\section*{Hoofdstuk \ref{ch:b1:stratigraphic-scale} --- De stratigrafische schaal}

\begin{solution}{exo:b1:stratigraphic-scale:1}
(a)~Chronostratigrafisch, een systeem. (b)~Lithostratigrafisch, een \omterm{def:b1:stratigraphic-scale:lithostratigraphy}{formatie}.
(c)~Geochronologisch, een tijdvak (haar gesteente is de Vroeg-Krijt-serie).
(d)~Biostratigrafisch, een \omterm{def:b1:stratigraphic-scale:biozone}{biozone}. (e)~Chronostratigrafisch, een \omterm{def:b1:stratigraphic-scale:chronostratigraphy}{etage}.
(f)~Geochronologisch, een tijdvak. (g)~Lithostratigrafisch, een \omterm{def:b1:stratigraphic-scale:lithostratigraphy}{lid}.
\end{solution}

\begin{solution}{exo:b1:stratigraphic-scale:2}
Paleozoïcum $538.8 - 251.902 = \qty{286.9}{\Myr}$, \qty{53.2}{\%} van het
Fanerozoïcum; Mesozoïcum $251.902 - 66.00 = \qty{185.9}{\Myr}$, \qty{34.5}{\%};
Cenozoïcum \qty{66.0}{\Myr}, \qty{12.2}{\%}. Het Fanerozoïcum is
$538.8/4567 = \qty{11.8}{\%}$ van de tijd sinds de basis van het Hadeïcum.
\end{solution}

\begin{solution}{exo:b1:stratigraphic-scale:3}
A, B, C, D werden in die volgorde afgezet (superpositie), horizontaal (oorspronkelijke horizontaliteit), dus is hun kanteling een latere gebeurtenis. De dijk E doorsnijdt A tot D, dus is zij jonger dan D (doorsnijding); of zij vóór of na het kantelen kwam, kan uit de beschrijving niet worden afgeleid. Het erosieoppervlak snijdt D en E af, dus is het jonger dan beide (doorsnijding). F en G werden daarna erop afgezet (superpositie). De breuk verplaatst A tot F maar niet G: zij bewoog na F en vóór G (doorsnijding). Volgorde: A, B, C, D; kanteling en intrusie van E (volgorde onbeslist); erosie; F; breukvorming; G.
\end{solution}

\begin{solution}{exo:b1:stratigraphic-scale:4}
$r = 30/(85.40 - 84.60) = \qty{37.5}{m}$ per miljoen jaar. Het eerste voorkomen
ligt $12/37.5 = \qty{0.32}{\Myr}$ na het onderste asbed: $85.40 - 0.32 =
\qty{85.08}{\Ma}$.
\end{solution}

\begin{solution}{exo:b1:stratigraphic-scale:5}
(a)~Vier verticale balken tegen de hoogte. (b)~B en C komen samen voor van
\qtyrange{26}{50}{m}: een zone van gelijktijdig bereik van \qty{24}{m} dik. De \omterm{def:b1:stratigraphic-scale:biozone}{intervalzone}
van het eerste voorkomen van B tot dat van D loopt van \qtyrange{12}{41}{m},
\qty{29}{m}. (c)~In de tweede sectie kunnen de eerste voorkomens van B (\qty{3}{m})
en van C (\qty{9}{m}) worden geplaatst, zodat de basis van de \omterm{def:b1:stratigraphic-scale:biozone}{intervalzone} van (b) en
de basis van de zone van gelijktijdig bereik worden herkend; haar top (het eerste
voorkomen van D) en de top van de overlap zoals in sectie~1 gezien, kunnen niet worden gecontroleerd, D ontbreekt.
De eerste voorkomens van B en C ($\qty{12}{m} \leftrightarrow
\qty{3}{m}$, $\qty{26}{m} \leftrightarrow \qty{9}{m}$) zijn de veiligste lijnen; de laatste
voorkomens (B op \qty{50}{m} en \qty{15}{m}) hangen meer af van de bewaring.
\end{solution}

\begin{solution}{exo:b1:stratigraphic-scale:6}
(a)~$60/20 = \qty{3}{\Myr}$: de top van het zandsteen is drie miljoen jaar
jonger in de landwaartse put. (b)~Één miljoen jaar slib van \qty{30}{m} per
miljoen jaar: het asbed ligt \qty{30}{m} boven de top van het zandsteen.
(c)~Toen de as viel, was de kustlijn \qty{20}{km} landinwaarts van de zeewaartse put verschoven, dus lag zij nog \qty{40}{km} zeewaarts van de landwaartse put, die op droog land stond. De as viel daar op het landoppervlak, onder het zandsteen dat de zee twee miljoen jaar later afzette --- als zij niet eerst werd weggespoeld, zoals een losse laag op het land heel goed kan. Hetzelfde asbed ligt in de ene put boven de formatiegrens en in de andere eronder: de grens snijdt de tijd.
\end{solution}

\begin{solution}{exo:b1:stratigraphic-scale:7}
(a)~Lochkovien $1.36/419.62 = \qty{0.32}{\%}$; Induan $0.024/251.902 =
\qty{0.0095}{\%}$, ongeveer 34 keer kleiner. De basis van het Induan ligt tussen asbedden die
precies zijn gedateerd in de GSSP-sectie zelf; een grote onzekerheid wijst op een grens
waarvan de leeftijd moet worden overgebracht vanaf verder gelegen gedateerde niveaus.
(b)~$251.941 - 251.880 = \qty{0.061}{\Myr} = 61$ duizend jaar, met de
onzekerheid $\sqrt{37^2 + 31^2} = 48$ duizend jaar: de gepubliceerde
\num[uncertainty-mode=separate]{60 \pm 48} duizend jaar.
\end{solution}

\begin{solution}{exo:b1:stratigraphic-scale:8}
(a)~$m = 3000 \times \frac43\pi\,(5000)^3 = \qty{1.57e15}{kg}$; het stof,
$60m = \qty{9.42e16}{kg}$. (b)~Het oppervlak van de aarde is $4\pi(6.378\times10^6)^2
= \qty{5.11e14}{m^2}$: \qty{184}{kg} per vierkante meter. (c)~$184/2500 =
\qty{0.074}{m}$, ongeveer \qty{7}{cm}. (d)~De dikte gaat als de derde macht van de
diameter: \qty{1.6}{cm} voor \qty{6}{km}, \qty{20}{cm} voor \qty{14}{km}. De orde
van grootte, centimeters tot decimeters, is die van de grensklei,
\qty{50}{cm} dik bij El Kef; de vergelijking kan alleen grof zijn, omdat het stof niet
gelijkmatig kan zijn gevallen en een klei ook het gewone sediment van haar tijd verzamelt.
\end{solution}

\begin{solution}{exo:b1:stratigraphic-scale:9}
(a)~$\num{60e3} \times \num{2e-6} = \qty{0.12}{m}$, \qty{12}{cm}; voor 12 en 108
duizend jaar, \qty{2.4}{cm} en \qty{21.6}{cm}. (b)~$\num{32e3} \times
\num{10e-6} = \qty{0.32}{m}$. (c)~Een extinctie beslaat centimeters tot decimeters
gesteente: monsters moeten centimeter voor centimeter worden genomen, en een enkel
monster van \qty{10}{cm} dik in de geconcentreerde sectie middelt vijftig duizend jaar ---
het grootste deel van de gebeurtenis.
\end{solution}

\begin{solution}{exo:b1:stratigraphic-scale:10}
(a)~$u = R/s = 10 \times 2000 = \qty{20000}{m}$ per miljoen jaar,
\qty{20}{km} per miljoen jaar; over \qty{30}{km}, $\Delta t = \qty{1.5}{\Myr}$.
(b)~De vlakte is nu het oppervlak $z = s x + a t$. De kustlijn voldoet aan
$s x + a t = z_0 + R t$, dus $x(t) = (z_0 + (R - a)t)/s$ en
$\mathrm dx/\mathrm dt = (R - a)/s$, positief (landinwaarts) zolang $a < R$. Wanneer
$a > R$ is de snelheid negatief: de kustlijn verschuift zeewaarts, een regressie, en
het kustzandsteen is dan zeewaarts jonger. (c)~Een klif van \qty{200}{m} lang beslaat
bij \qty{20}{km} per miljoen jaar $0.2/20 = \qty{0.01}{\Myr}$ --- tienduizend
jaar diachronie, ver onder wat fossielen of de meeste dateringen kunnen oplossen; het gesteente
aan weerszijden van het contact is overal langs het contact hetzelfde.
\end{solution}

\begin{solution}{exo:b1:stratigraphic-scale:11}
(a)~$f = (25 - 10)/(50 - 10) = 0.375$ en $t = 100.50 - 0.375 \times 0.80 =
\qty{100.20}{\Ma}$, met
\[
  \sigma_t = \sqrt{(0.625 \times 0.05)^2 + (0.375 \times 0.03)^2} = \qty{0.033}{\Myr}.
\] (b)~$g(f) = (1 - f)^2\sigma_1^2 + f^2\sigma_2^2$
heeft $g'(f) = -2(1 - f)\sigma_1^2 + 2f\sigma_2^2$, nul bij $f^* =
\sigma_1^2/(\sigma_1^2 + \sigma_2^2) = 0.0025/0.0034 = 0.735$, dichter bij het beter
gedateerde asbed; daar is $\sigma_t = \sigma_1\sigma_2/\sqrt{\sigma_1^2 + \sigma_2^2} =
\qty{0.026}{\Myr}$, kleiner dan beide. (c)~$g$ is een parabool die naar boven opent, dus
op $[0, 1]$ is zij het grootst aan een uiteinde: $g(0) = \sigma_1^2$, $g(1) = \sigma_2^2$.
Dus $\sigma_t \le \max(\sigma_1, \sigma_2)$ tussen de lagen.
\end{solution}

\begin{solution}{exo:b1:stratigraphic-scale:12}
(a)~De grens ligt halverwege in diepte: $f = 0.5$, \qty{29.50}{\Ma}. (b)~Met
$\varphi_0\lambda = 0.63 \times 1960 = \qty{1234.8}{m}$ is de vaste dikte van
\qtyrange{600}{1000}{m} $400 - 1234.8\,(e^{-600/1960} - e^{-1000/1960}) =
\qty{232.2}{m}$, en van \qtyrange{200}{600}{m}, \qty{194.2}{m}; samen
\qty{426.3}{m}. Het aandeel van de korrels dat is afgezet tussen het onderste asbed en de grens is
$232.2/426.3 = 0.545$, en de leeftijd is $30.0 - 0.545 =
\qty{29.46}{\Ma}$. (c)~De tweede schatting is \num{45} duizend jaar jonger.
Zij is de betere: de diepere helft heeft meer van haar water verloren, dus bevat
zij per meter meer korrels, en dus meer tijd, dan de ondiepere helft (gemiddelde
porositeit \num{0.42} tegenover \num{0.51}).
\end{solution}

\subsection*{Probleem: drie secties over een bekken correleren}

\begin{solution}{pb:b1:stratigraphic-scale:1}
\begin{enumerate}
  \item Drie kolommen op, zeg, \qty{1}{cm} voor \qty{10}{m}, verbonden door de asbedden en door de zonegrenzen; de wig van de Kell Sandstone wordt van \qty{45}{m} in A tot \qty{6}{m} in C dunner.
  \item Zone~2: A $65.8 - 36.6 = \qty{29.2}{m}$, B \qty{20.8}{m}, C \qty{10.0}{m}. Zone~3 gelogd: A \qty{14.2}{m}, B \qty{14.0}{m}, C \qty{7.2}{m}.
  \item Zone~1 begint bij het eerste voorkomen van \emph{V.~primum}, dat vanaf de basis van elke sectie wordt gevonden: het eerste voorkomen ligt eronder.
  \item Hij leeft alleen in het zandsteen: hij volgt de facies, en de facies zijn van plaats tot plaats van verschillende leeftijd. In A komt hij voor in zones~1 en~2, in B en C alleen in zone~1.
  \item In zone~2 in alle drie (A \qty{49.0}{m}, B \qty{40.0}{m}, C \qty{19.0}{m}). Hij is hier van sectie tot sectie consistent, en zo zou hij zone~2 kunnen onderverdelen, maar een laatste voorkomen hangt af van de bewaring; het moet eerst in meer secties worden gecontroleerd.
  \item De lijnen verbinden \qty{36.6}{m}, \qty{31.2}{m} en \qty{14.8}{m}, dan \qty{65.8}{m}, \qty{52.0}{m} en \qty{24.8}{m}. Zone~2 dunt uit in de verhouding $29.2 : 20.8 : 10.0$, ongeveer $2.9 : 2.1 : 1$: als de zone overal even lang duurde, is het sediment in A ongeveer drie keer zo snel opgehoopt als in C.
  \item A: zone~2 ($45.0 > 36.6$); B: zone~1 ($20.0 < 31.2$); C: zone~1 ($6.0 < 14.8$).
  \item Nee: hij ligt in zone~2 in A, zone~1 in B en C, dus is hij in het westen jonger. De mariene klei breidde zich westwaarts uit over het kustzand: de kustlijn verschoof westwaarts, landinwaarts, in een transgressie.
  \item Een as valt in een geologisch ogenblik over het hele bekken en bezinkt op elk sediment; de top van het zandsteen volgt de bewegende kustlijn en is \omterm{def:b1:stratigraphic-scale:diachronous}{diachroon}.
  \item Laterale continuïteit (de laag loopt over het bekken), en superpositie (as~1 onder as~2 in elke sectie, wat de correlatie controleert).
  \item A $(71.0 - 5.0)/1.20 = \qty{55.0}{m}$ per miljoen jaar; B $48.0/1.20 = \qty{40.0}{}$; C $24.0/1.20 = \qty{20.0}{}$.
  \item A ligt het dichtst bij het land, waar het sediment werd aangevoerd; minder ervan bereikte het diepere water naar het oosten.
  \item $0.1/55 = \num{1.8}$ duizend jaar in A, $\num{2.5}$ in B, $\num{5.0}$ in C: A lost de tijd het fijnst op.
  \item A $152.40 - 40.0/55.0 = \qty{151.67}{\Ma}$; B $152.40 - 12.0/40.0 =
        \qty{152.10}{\Ma}$; C $152.40 - 3.0/20.0 = \qty{152.25}{\Ma}$.
  \item C naar B: $20/(152.25 - 152.10) = \qty{133}{km}$ per miljoen jaar; B naar A:
        $25/(152.10 - 151.67) = \qty{58.5}{km}$ per miljoen jaar. De transgressie vertraagde.
  \item A $152.40 - 60.8/55.0 = \qty{151.29}{\Ma}$; B $152.40 - 44.0/40.0 =
        \qty{151.30}{\Ma}$; C $152.40 - 21.8/20.0 = \qty{151.31}{\Ma}$.
  \item Zij beslaan \num{16} duizend jaar, minder dan de onzekerheid van elk (ongeveer \num{28} duizend jaar, vraag~19): het eerste voorkomen is over het bekken synchroon binnen de resolutie van de dateringen, wat het een goede indexwaarde geeft.
  \item A $152.40 - 31.6/55.0 = \qty{151.826}{\Ma}$; B $152.40 - 23.2/40.0 =
        \qty{151.820}{\Ma}$; C $152.40 - 11.8/20.0 = \qty{151.810}{\Ma}$.
  \item $f = 0.479$, $0.483$, $0.492$; $\sigma_t = \sqrt{(1 - f)^2 + f^2} \times
        0.04 = \qty{0.028}{\Myr}$ in elk, \num{28} duizend jaar.
  \item Gemiddelde \qty{151.819}{\Ma}; spreiding $151.826 - 151.810 = \num{16}$ duizend
        jaar, kleiner dan de onzekerheid van elke leeftijd: de drie komen overeen.
  \item De interpolatie gaat uit van één snelheid van as~1 tot as~2. In A bevat dat interval kustzandsteen tot \qty{45}{m} en daarboven schalie; zand en slib hoeven niet met dezelfde snelheid te zijn afgezet, en zij compacteren verschillend (\cref{ch:b1:diagenesis-basins}). Als het zand sneller is afgezet, bevat het zandsteendeel minder tijd dan zijn dikte suggereert, en is de eerste verschijning ouder dan berekend --- de leeftijd uit A is inderdaad de oudste van de drie.
  \item B: het eerste voorkomen ligt in mariene schalie, \qty{11}{m} boven de top van het zandsteen, zonder verandering van facies nabij het niveau, en de snelheid geeft \num{2.5} duizend jaar per \qty{10}{cm}. A lost de tijd beter op, maar plaatst het niveau in het kustzandsteen, \qty{8.4}{m} onder een verandering van facies; C is geconcentreerd, \num{5} duizend jaar per \qty{10}{cm}. C is het slechtst voor de resolutie, A voor de facies.
  \item Uit de gemiddelde leeftijden van de twee eerste voorkomens: $151.819 - 151.302 =
        \qty{0.52}{\Myr}$: ongeveer een halve miljoen jaar.
  \item De grens is \qty[uncertainty-mode=separate]{151.82(3)}{\Ma}.
\end{enumerate}
\end{solution}
